% ============================================================================
% Chapter 6 · Conclusion
% Style: no em dashes; synthesise, do not re-argue; answer each RQ directly;
% state each headline number once; limitations and future work relocated here
% from the discussion to avoid duplication.
% ============================================================================

\section{Conclusion}\label{sec:conclusion}

In 2024 the German day-ahead price fell below zero for 459 hours, and this
thesis set out to establish what happens to the value of offshore wind once
hours like these become a sea basin's ordinary operating state. Europe is
committed to the buildout that will produce them: Baltic offshore capacity is
planned to grow from 2.6 to 49.5~GW by 2040, a nineteenfold increase against a
ninefold increase around the North Sea, while support moves away from fixed
payments towards contracts that leave generators exposed to the wholesale
price. What a farm earns therefore depends on the price prevailing in the
hours it generates, and those are the hours its neighbours generate too.

This thesis measured the share of the average wholesale price that offshore
wind actually receives, its capture value, for every bidding zone around the
two seas, and projected it to 2030 and 2040. Over 2023 and 2024 the Baltic
zones that have built little so far captured 0.883 of the average price, level
with the North Sea's 0.879. By 2040, on today's grid, they fall to 0.47 under
National Trends, the pathway of the official plan that tracks existing
national policies, while the North Sea holds 0.57, so close to half of present
value is lost. The ordering of the two basins reverses within fifteen years.

The largest single cause at the level of an individual farm is crowding,
meaning how closely a farm generates at the same time as everything else, and
the penalty for it is twice as steep in the Baltic as in the North Sea. At the
level of the system the fall divides in two. Of the 0.42 the Baltic loses,
0.23 comes from the buildout itself and 0.19 from congestion, and the
congestion part disappears once the grid is expanded: with transmission free
to grow the Baltic recovers to 0.66 against the North Sea's 0.64. Nearly half
the loss, in other words, is a consequence of what gets built around the
farms.

Which of the two models could make that projection had to be established
rather than assumed. Over the years that can be checked against observation
the statistical model and the system model are equally accurate, and their
estimates of capture value agree to within 0.005. Beyond those years ninety
per cent of their disagreement at 2040 is the statistical model ceasing to
respond to the scenario, because it inherits the present instead of
reconstructing it. The statistical model therefore measures the present and
the near term, and the system model, which rebuilds prices from plants, costs
and the network, carries the projection.

The first of this work's four contributions is a systematic comparison of
offshore capture value between the North Sea and the Baltic at the resolution
of the bidding zones in which prices are set, carried from today's measurement
through to the 2040 buildout. The second sets a statistical model and a system
model against observation and against each other on the same measured
quantity, and establishes by measurement the point at which a projection must
pass from one to the other. The third divides the loss of capture value twice,
into the part no network removes and the congestion part an expanded network
does, and into the shares of the loss borne by the farm and by the state under
the contracts that now dominate support. The fourth computes a value-adjusted
cost of energy for each location rather than for each technology, so that a
measured fall in value can be read as the price a specific site must earn.

What the Baltic buildout earns will be decided by the investments delivered
alongside it. Expanding transmission removes most of the congestion loss, and
the income the congested grid already generates, which European law sets aside
for exactly this use, would cover the planned transmission programme several
times over. Allowing demand to move into the windy hours is worth as much
again. The value of Baltic offshore wind is therefore set by decisions that
are still open, and most of that capacity has not yet been built.